% Section 3.2 -- instrument characterization, part one.
% Companion to the dag_layer finding. Framing: a caution for consumers of the
% MathlibGraph release, not a limitation of this work. No literal numbers:
% every figure is a macro from numbers.tex. The retention decomposition lives
% in 3.3 -- this section states the set relation and points forward.

\subsection{\texttt{is\_explicit} does not mean ``typed in source''}
\label{sec:is-explicit}

The dependency graph we build on~\citep{li2026network} distinguishes \emph{explicit} edges from the
much larger set of compiler-synthesized ones, and identifies the explicit
subgraph as the closer proxy for human-intended dependency structure. That
subgraph is our object of study, and the distinction is the right one to draw.
The flag's name, however, invites a reading it does not support, and anyone
building a source-level pipeline on this release should know which reading is
correct before choosing a denominator.

\paragraph{Where the flag comes from.}
We consume the published graph rather than re-running its extraction; what
follows is provenance, read off the release's own parser and the documented
behaviour of the tool upstream of it. \texttt{is\_explicit} originates as the
\texttt{*} marker emitted by \texttt{lean-training-data}'s \texttt{premises}
command~\citep{leantrainingdata}, and the release's parser consumes it with that
meaning. The marker records that a constant occupies an \emph{explicit argument
position of the elaborated term} (a position not filled by instance
resolution or unification). It is a property of the term the kernel accepts, and
by construction not a claim about the surface syntax a human wrote, since it is
assigned after elaboration, when the source text is gone.

\paragraph{The two readings come apart at the top of the graph.}
Elaboration inserts constants of its own. In the published graph,
proof-irrelevance and rewriting plumbing (\texttt{propext}, \texttt{Eq.mpr},
\texttt{of\_eq\_true}) carry explicit in-edges throughout the explicit
subgraph and are among its most heavily cited nodes by in-degree. No mathlib
author writes any of them. They are artifacts of how \texttt{rw}, \texttt{simp}
and the equation compiler construct terms, and they are explicit in the sense
the flag means, sitting in explicit argument positions of terms that no human
typed in that form. The same mechanism operates less visibly throughout
the library, in notation-backed definitions that are never written as
identifiers, in constants emitted by arithmetic normalization, and in simp-set
lemmas that fire without being named.

\paragraph{The consequence for any source-level instrument.}
We take the narrower of two available claims, and it is enough. Some explicit
dependencies are never written as identifiers at all (the plumbing above is the
clearest case), so no instrument that reads source can recover them, and
\emph{no} source-level instrument, including a hypothetically perfect one with
exact name resolution, complete namespace and \texttt{open} handling and no
tokenization loss, should be expected to reach unit retention against the
explicit set. We do not claim the stronger set relation. The two sets are not
nested in either direction: a name written in source need not survive as the
same dependency in the elaborated term, and source text carries names that are
not cited constants at all. This is why we report retention against a scope
decomposition rather than against a threshold.
Section~\ref{sec:retention} gives that decomposition as three quantities cut
with one matcher at three scopes, and says what each does and does not
license.

Two consequences follow for anyone building on this release. High retention
rates reported for source-visible extraction elsewhere~\citep{cslibpremise} are
not comparable to ours, because a retention figure inherits the semantics of the
artifact it is scored against and belongs to an (instrument, ground-truth) pair
rather than to the instrument alone. And membership in the explicit set is not
evidence that a human wrote the name, so it should not be used to build a
retrieval target, to score a premise-selection model against ``the premises a
mathematician used,'' or to set a parser's accuracy bar. A check that costs
very little catches this, and it is the same one that catches the attribute of
Section~\ref{sec:dag-layer}: take the highest-in-degree nodes and ask whether a
mathlib author would type them.
